% ================= 第 3 章 (后半): 3.3 猜想陈述 =================

\subsection{陈述引言中的局部-整体相容性猜想 1.1.1}\label{sec:3.3}

固定一个如 \ref{sec:3.2} 的模形式, 其中 $f$ 在 $\Gamma_1(Mp^{r})$ 上为新, 并进一步假设 $\pi_p(f)$ 是主级数, 使得可以写成:
\[
(\Sym^{k-2}L^{2})^{\vee} \otimes_L \pi_p(f) = \pind{\mathrm{lp},\le k-2}_{B(\Qp)}^{\GL_2(\Qp)} \chi_1 \otimes \chi_2,
\]
其中 $\chi_1 \otimes \chi_2 \in \widehat{T}(L)$ 是经典权 $k$ 的, $\chi_2$ 局部常值, 且 $\chi_1\chi_2 = \varepsilon^{2-k}\chi_p^{-1}$ ($\chi_p$ 是把 $\chi$ 看作 $\mathbb{A}_{f}^{\times}$ 的特征标时其在 $p$ 处的局部分量).
设 $\sigma_p(f) := \sigma(f)|_{\Gal(\overline{\Qp}/\Qp)}$.
在这种情况下, 由 \cite{ref:40}, $\sigma_p(f)$ 在完全分歧的 Abel 扩张 $\Qp(\zeta_{p^{r}})$ 上变成晶体.
最后假设 $\sigma_p(f)$ 可约, 或者等价地 (必要时用光滑诱导上的缠结修改 $\chi_1$ 和 $\chi_2$), $\val(\chi_2(p)) = k-1$ 和 $\val(\chi_1(p)) = 1-k$.
设 $\eta_1 := \chi_1|\,|^{1-k}\varepsilon^{k-2}$, $\eta_2 := \chi_2|\,|^{k-1}\varepsilon^{2-k}$, 且如 \ref{sec:2.3} 有:
\[
\sigma_p(f) \simeq \begin{pmatrix} \chi_1|\,|^{1-k}\varepsilon^{k-2} & * \\ 0 & \chi_2|\,|^{k-1}\varepsilon^{1-k} \end{pmatrix} = \begin{pmatrix} \eta_1 & * \\ 0 & \eta_2\varepsilon^{-1} \end{pmatrix}.
\]
$*$ 是否为零 (回顾: 若 $*$ 非零, 则只存在一个可能的扩张) 与滤过 $\phi$-模 $D_p(f) := D_{\mathrm{cris}}(\sigma_p(f))$ 上 Hodge 滤过的参数有关.
事实上容易证明 (参见 \cite{ref:31}):
\[
\begin{aligned}
&D_p(f) = Le_1 \oplus Le_2,\\
&\phi(e_1) = \chi_2(p)e_1,\\
&\phi(e_2) = p^{k-1}\chi_1(p)e_2,\\
&\mathrm{Fil}^{i}D_p(f)_{\Qp(\zeta_{p^{r}})} = D_p(f)_{\Qp(\zeta_{p^{r}})}, \qquad i \le 0,\\
&\mathrm{Fil}^{i}D_p(f)_{\Qp(\zeta_{p^{r}})} = \Qp(\zeta_{p^{r}}) \otimes_{\Qp} L \cdot (e_1 + \delta x e_2), \qquad 1 \le i \le k-1,\\
&\mathrm{Fil}^{i}D_p(f)_{\Qp(\zeta_{p^{r}})} = 0, \qquad i \ge k,
\end{aligned}
\]
其中 $D_p(f)_{\Qp(\zeta_{p^{r}})} := \Qp(\zeta_{p^{r}}) \otimes_{\Qp} D_p(f)$, $x \in (\Qp(\zeta_{p^{r}}) \otimes_{\Qp} L)^{\times}$ 只依赖于特征标 $\chi_1z^{k-2}$ 和 $\chi_2$, 而 $\delta \in \{0,1\}$ 是 Hodge 滤过的"参数".
还存在 $\Gal(\Qp(\zeta_{p^{r}})/\Qp)$ 在 $D_p(f)$ 上保持 $\mathrm{Fil}^{k-1}D_p(f)_{\Qp(\zeta_{p^{r}})}$ 的作用, 这里不给出 (参见 \cite{ref:31}).
可见 $\sigma_p(f)$ 分裂当且仅当 $\delta = 0$.

\begin{proposition}\label{prop:3.3.1}
有拓扑 $\GL_2(\Qp)$-等变同构:
\[
\widehat{\pi}_p(f) \simeq \cind_{B(\Qp)}^{\GL_2(\Qp)} \eta_2 \otimes \eta_1.
\]
特别地, $\widehat{\pi}_p(f)$ 拓扑不可约, 并等同于 $(\Sym^{k-2}L^{2})^{\vee} \otimes \pi_p(f)$ 关于 $\mathcal{O}_L[\GL_2(\Qp)]$ 上某个 (任意) 有限型 $\mathcal{O}_L$-格的完备化.
\end{proposition}

\begin{proof}
命题陈述的第二部分, 即 $(\Sym^{k-2}L^{2})^{\vee} \otimes_L \pi_p(f)$ 关于一个 $\GL_2(\Qp)$-稳定的有限型 $\mathcal{O}_L$-格的完备化的描述及其性质, 是命题 \ref{prop:2.2.1} 的内容.
第一部分由如下事实推出: $(\Sym^{k-2}L^{2})^{\vee} \otimes_L \pi_p(f)$ 中所有 $\GL_2(\Qp)$-稳定的 $\mathcal{O}_L$-格, 特别是与 $\widehat{H}{}_{c}^{1}(K_1^{p}(M))$ 相交诱导的那个, 都给出等价范数, 这容易从第二部分推出 (参见 e.g.\ \cite[Cor.~5.4.4]{ref:3}).
\end{proof}

另一方面, 我们在 \ref{sec:2.3} 给可约表示 $\sigma_p(f)$ 联系了一个长度为 $2$ 的酉 $\GL_2(\Qp)$-Banach 空间 $B(\sigma_p(f))$, 它是 $\cind_{B(\Qp)}^{\GL_2(\Qp)} \eta_1\varepsilon \otimes \eta_2\varepsilon^{-1}$ 被 $\widehat{\pi}_p(f)$ 的扩张, 且 (作为 $\GL_2(\Qp)$ 的拓扑表示) 分裂当且仅当 $\sigma_p(f)$ (作为 $\Gal(\overline{\Qp}/\Qp)$ 的表示) 分裂.

\begin{conjecture}\label{conj:3.3.2}
闭浸入 $\widehat{\pi}_p(f) \hookrightarrow \Pi_p(f)$ 可延拓为拓扑 $\GL_2(\Qp)$-等变同构:
\[
B(\sigma_p(f)) \xrightarrow{\sim} \Pi_p(f).
\]
\end{conjecture}

这个猜想 (我们在后面证明它的一个弱版本) 应被看作"$p$ 进 Langlands 对应"框架下的一种局部-整体相容性.
事实上, 它断言: 整体理论诱导的 $\GL_2(\Qp)$ 表示 (表示 $\Pi_p(f)$) 是一个"Fontaine 意义下"的局部表示 (表示 $B(\sigma_p(f))$).
如同半稳定情形 (\cite{ref:8}, \cite{ref:9}, \cite{ref:18}), Banach 空间 $\widehat{H}{}_{c}^{1}(K_1^{p}(M))_L$ 的 $\sigma(f)$-等型分量应当因此包含与 $\sigma_p(f)$ 或其滤过模的给出精确"等价"的信息.

当 $\sigma_p(f)$ 分裂时, 我们将看到出现在 $\Pi_p(f)$ 中的第二个块, 即 $\GL_2(\Qp)$-Banach 空间 $\cind_{B(\Qp)}^{\GL_2(\Qp)} \eta_1\varepsilon \otimes \eta_2\varepsilon^{-1}$, 是这种情况下来自经典形式 $f$ 的超收敛伴随模形式存在的表现.
它的模 $p$ 约化恰好是 Serre 意义下 $f$ 的伴随形式 (\cite{ref:33}, \cite{ref:37}, \cite{ref:29}).
